\documentclass[11pt,a4paper]{article}
\usepackage[T1]{fontenc}
\usepackage{lmodern,amsmath,amssymb,booktabs,microtype,array}
\usepackage[margin=24mm]{geometry}
\usepackage[colorlinks=true,urlcolor=blue,linkcolor=blue]{hyperref}
\newcommand{\HH}{\mathcal H}
\newcommand{\Dom}{\operatorname{Dom}}
\newcommand{\Sp}{\operatorname{Sp}}
\title{Domains, phase singularities, and the\\two-variable pentagon benchmark}
\author{Alan Day \and Chengheng\thanks{AI research assistant (OpenAI).}}
\date{Working research note --- 1 November 2026}
\begin{document}\maketitle
\begin{abstract}
We revisit an unresolved reproduction of the positive-cosmological-constant
local phase operator. Before further computation, we trace the existing
proofs and their hypotheses. In a specified Schr\"odinger representation,
we identify a strong-commutation obstruction, singular domains of
individual chiral factors, and loss of first-order Sobolev regularity
under the bounded paired multiplier. These elementary domain diagnostics
do not refute the bounded pentagon. Eleven oversampled Hermite calculations
are preserved with their projection losses; their quadrature error remains
uncertified. No finite-$G$ observer gluing or holographic dual is constructed.
\end{abstract}
\section{Objects, state, and scope}
All coordinates are dimensionless, $\hbar>0$, and
$\HH=L^2(\mathbb R^2,dx\,dy)$ with Lebesgue measure. Let
\[
 X=x,\quad Y=y,\quad P_x=-\pi i\partial_x,\quad P_y=-\pi i\partial_y,
 \qquad \psi(x,y)=\pi^{-1}e^{-(x^2+y^2)/(2\pi)}.
\]
The self-adjoint position and momentum closures obey their Weyl relations;
$\|\psi\|=1$. Use the existing scalar functions \cite{ks,kim}
\begin{align}
 \psi_q(w)&=\prod_{n\geq0}(1+q^{2n+1}w)^{-1},\\
 \Phi^H(z)&=\frac{\psi_{e^{\pi iH}}(e^z)}
 {\psi_{e^{-\pi i/H}}(e^{z/H})},\qquad \Im H>0,\\
 F_\hbar(x,y)&=\frac{\Phi^{i\hbar}(x+i\hbar y)}
 {\overline{\Phi^{i\hbar}(x+i\hbar y)}}.
\end{align}
Assign arbitrary unit values at the discrete zeros and poles. Then
$A=F_\hbar(X,Y)$ is unitary by measurable functional calculus. Define
\[
 B=F_\hbar(P_y,P_x),\qquad M=e^{iXY/\pi},\qquad
 C=M^{-1}BM=F_\hbar(X+P_y,Y+P_x).
\]
Our unresolved test is $AB=BCA=BM^{-1}BMA$. This is the separately
stated local identity, not the full geometric flip with all linear factors.
The representation is a mathematical benchmark: $\hbar$ is not here
identified with an observer mass or a derived finite-$G$ parameter.

\newpage
\section{Formal commutation and the bounded phase}
Consider the normal operators
\[
 \overline Z_1=X-i\hbar Y,\qquad Z_2=P_y+i\hbar P_x.
\]
Their algebraic commutator vanishes on the Schwartz space. Nevertheless,
their spectral measures do not commute. Their spectral real and
imaginary parts include $X$ and $\hbar P_x$, for which
\[
 e^{iaX}e^{ib\hbar P_x}
   =e^{-i\pi\hbar ab}e^{ib\hbar P_x}e^{iaX}.
\]
Taking $a=1$, $b=1/\hbar$ gives phase $-1$ and a commutator of norm $2$
on every unit vector. Strong commutation of the two normal operators
would imply commutation of these bounded spectral functions, a contradiction.
This tests the literal joint-spectral interpretation of the cross-chiral
step in KS, \S6.4. It does not invalidate commutation of $Z_1$ with its
\emph{own} adjoint, nor rule out a different domain-based identity.

\subsection*{A bounded multiplier need not preserve a smooth core}
The first simple pole of $\Phi^{i\hbar}$ is
\[
 z_0=\pi\hbar-\pi i,\qquad
 (x_0,y_0)=(\pi\hbar,-\pi/\hbar),\qquad
 \zeta=(x-x_0)+i\hbar(y-y_0).
\]
Write $\Phi^{i\hbar}(z_0+\zeta)=R\,u(\zeta)/\zeta$,
where $R\ne0$, $u$ is analytic, and $u(0)=1$.
Thus, off the pole,
\[
 F_\hbar(x,y)=\frac{R}{\overline R}
       \frac{\overline\zeta}{\zeta}
       \frac{u(\zeta)}{\overline{u(\zeta)}}.
\]
The smooth last factor cannot remove the winding $-2$.
For $\zeta=r e^{i\theta}$ the leading squared gradient is
$4(\sin^2\theta+\hbar^2\cos^2\theta)/r^2$.
Since $dx\,dy=r\,dr\,d\theta/\hbar$, it follows that
\begin{align}
 \int_{\epsilon<|\zeta|<\rho}|\nabla(A\psi)|^2\,dx\,dy
 &=D_\hbar\log(\rho/\epsilon)+O(1),\quad \epsilon\downarrow0,\\
 D_\hbar&=\frac{4(1+\hbar^2)}{\pi\hbar}
       e^{-\pi(\hbar^2+\hbar^{-2})}>0 .
\end{align}
Here $\rho$ is fixed and small enough to exclude other zeros/poles.
The derivatives away from the pole must agree with any weak derivatives;
their divergent norm proves $A\psi\notin H^1(\mathbb R^2)$.
Each partial derivative separately diverges. In particular,
$A\psi\notin\Dom P_x,\Dom P_y$.
Unitarity on $L^2$ survives. This does not exclude a singular adapted core,
and by itself says nothing about equality of two bounded operator words.

\newpage
\section{Individual factors and moving domains}
The first pole of $\Phi^H$, $z_0(H)=-\pi i(1+H)$, has residue
\[
 R_H=-\frac{(\widetilde q^2;\widetilde q^2)_\infty}
                {(q^2;q^2)_\infty},
 \qquad q=e^{\pi iH},\quad\widetilde q=e^{-\pi i/H},
 \quad (a;q)_\infty=\prod_{n\geq0}(1-aq^n).
\]
For $\Im H>0$ both products converge and the residue is nonzero.
This follows by separating the first vanishing denominator
factor in the known product; it is not a new special-function identity.
For $H=i\hbar$ the Gaussian is nonzero at $(x_0,y_0)$, and
\begin{align}
 \int_{\epsilon<|\zeta|<\rho}
       |\Phi^{i\hbar}(x+i\hbar y)\psi(x,y)|^2\,dx\,dy
 &=C_\hbar\log(\rho/\epsilon)+O(1),\\
 C_\hbar&=\frac{2|R_{i\hbar}|^2}{\pi\hbar}
       e^{-\pi(\hbar^2+\hbar^{-2})}>0 .
\end{align}
Consequently $\psi\notin\Dom\Phi^{i\hbar}(X+i\hbar Y)$.
At $\hbar=1$, $R_i=-1$ and $C_1=2e^{-2\pi}/\pi
\simeq0.00118885096677$; $D_1=4C_1$.
The bounded ratio cannot automatically be factored into an everywhere
defined product of the two unbounded factors.

\subsection*{No common dense multiplication domain on a parameter open set}
For $H=a+ib$, $a,b>0$, the first pole meets the real coordinate plane at
\[
 x_0=\frac{\pi(a^2+a+b^2)}b,\qquad
 y_0=-\frac{\pi(a+1)}b.
\]
This maps the first quadrant diffeomorphically onto
$W=\{x>0,y<0,xy<-\pi^2\}$. The inverse is
\[
 a=\frac{-xy-\pi^2}{\pi^2+y^2},\qquad
 b=\frac{\pi(x-y)}{\pi^2+y^2}.
\]
For any nonempty open parameter set $\Omega$ in that quadrant, let
$W_\Omega$ be its image. Then
\[
 \bigcap_{H\in\Omega}\Dom\Phi^H(X+HY)
 \ \subset\ \{f\in L^2:f=0\ \text{a.e. on }W_\Omega\}.
\]
For arbitrary $L^2$ vectors put $g=|f|^2$.
At a Lebesgue point $p$ with $g(p)>0$, the ball mass
$G_p(r)=\int_{B_r(p)}g$ satisfies $G_p(r)\sim\pi g(p)r^2$.
Tonelli gives
\[
 \int_{B_\rho(p)}\frac{g(u)}{|u-p|^2}\,du
 =\frac{G_p(\rho)}{\rho^2}
       +2\int_0^\rho\frac{G_p(r)}{r^3}\,dr=\infty.
\]
The corresponding simple pole has this locally equivalent weight.
Thus $g$ must vanish almost everywhere on $W_\Omega$.
The right-hand side is a proper closed subspace, so the intersection is
not dense. This concerns the literal multiplication family on one fixed
Hilbert space; weak distributions, moving contours, parameter-dependent
vectors, and the paired bounded family are distinct.

\newpage
\section{Which existing proofs can be reused?}
We separate a theorem's published conclusion from applicability to this
realization. This is a local reading audit, not a claim
to have independently verified every argument in the cited papers.
\begin{center}\small
\begin{tabular}{p{.27\linewidth}p{.63\linewidth}}
\toprule
Source and locator & Conditions and dependencies\\
\midrule
Goncharov \cite{gon}, \S2.3--2.5 &
Positive real parameter; $L^2(\mathbb R)$; algebra-adapted test space
and topology. A complex two-variable extension is extra work.\\[4pt]
Fock--Goncharov \cite{fg}, Thm.~5.6, eq.~(87), Prop.~5.8 &
Contour displacement uses a pole-free analytic strip and rapid decay.
Extension uses the seminorms $\|Bf\|$, not only the Hilbert norm.\\[4pt]
Kim--Scarinci \cite{ks}, Prop.~6.11 &
Check cancellation and fixed-vector parameter continuation:
the literal single-factor family has the obstruction above.\\[4pt]
Kim \cite{pseudo}, Props.~7.3--7.4, 9.1--9.2 &
Prop.~9.1 uses Prop.~7.4; \S7 attributes that identity to KS.
This is not independent validation of the continuation step.\\[4pt]
Pusz--So{\l}tan \cite{ps}, Def.~2.1, Prop.~3.1 &
For $0<q<1$, regular pairs have spectra in
$\overline\Gamma$, $\Gamma=\{z:|z|\in q^{\mathbb Z}\}$, and specified
bicharacter covariance. Formal relations alone are insufficient.\\
\bottomrule
\end{tabular}
\end{center}
There is a concrete mismatch with the last route.
In our Lebesgue representation, $e^{X+i\hbar Y}$
has spectrum $\mathbb C$: its essential range is all nonzero complex
numbers, and its closure includes zero. This is not contained in
$\overline\Gamma$, a union of discrete-radius circles and zero.
That regular-pair theorem cannot simply be applied here.
It may motivate another representation, but such a change would need
its own geometric and physical justification.

The practical lesson is to use the established real-parameter proof
as a source of precise hypotheses. Scalar meromorphic identities,
equations on an analytic core, weak continuation of matrix
elements, and equality of bounded spectral operators are different claims.
The preceding domain tests exclude some naive identifications; they do
not settle all formulations of the paired identity.

Source versions are pinned in the bibliography and source record.
We inspected the full relevant rendered pages for FG eq.~(87), Kim's
use of Prop.~7.4, and the regular-pair definition.
Further candidates remain candidates until their assumptions and
normalization are matched; absence of a located correction is not
evidence that no correction exists.

\newpage
\section{Oversampled projection experiment}
Let $P_n$ project onto the first $n$ Hermite functions in each coordinate.
For an \emph{exact} compression $v_j=P_nU_jv_{j-1}$ of a unitary $U_j$,
\[
 \ell_j=\|(1-P_n)U_jv_{j-1}\|
       =\sqrt{\|v_{j-1}\|^2-\|v_j\|^2}.
\]
Telescoping gives an error at most $\sum_j\ell_j$ for a word.
Our left word uses $[B,A]$ and the right uses
$[A,M,B,M^{-1},B]$, in order of application.
Their projected difference minus the seven losses would be a continuous
lower bound if all compressions were exact.

We instead use $q>n$ Gauss--Hermite nodes. If $Q$ diagonalizes the $q$-node
Jacobi matrix and $W$ is its first $n$ rows, a multiplier step is
\[
 v\longmapsto W\bigl[f\odot(W^TvW)\bigr]W^T.
\]
The Fourier transform is diagonal with coefficients $(-i)^{m+k}$.
These are oversampled approximations, not exact Galerkin compressions.
Define $\mu_q=r_q-\sum_j\ell_{j,q}$.
\emph{Positive $\mu_q$ is not a certified continuous lower bound.}

\begin{center}\small
\begin{tabular}{rrrrr}
\toprule
$\hbar$ & $n$ & $q$ & residual $r_q$ & margin $\mu_q$\\
\midrule
1 &16&128& .106990701&-.101943670\\
1 &32&256& .114715070&-.019846987\\
1 &64&256& .118041741& .027620827\\
1 &64&512& .117989373& .026793971\\
1 &128&512& .119750912& .057061713\\
1 &128&768& .119712207& .056877670\\
1 &128&1024& .119630925& .057013742\\
.5&64&512& .028263351& .007989570\\
.5&128&768& .028962160& .016228235\\
flat&32&256& $4.6300\,10^{-5}$& $-5.0987\,10^{-5}$\\
flat&64&512& $2.5202\,10^{-7}$& $-4.9643\,10^{-7}$\\
\bottomrule
\end{tabular}
\end{center}
The flat control uses $F_0(x,y)=\exp[-iy\log(1+e^x)/\pi]$.
Compression loses norm even for this control; unlike nodewise DVR,
it need not be unitary in the finite subspace.

All intermediate coefficient arrays are saved, with input hashes and
Python/NumPy/SciPy versions. Independent saved-vector arithmetic replay
passes 322 checks. It does not re-evaluate the multiplier integrals.
Product tails, projection loss, quadrature, and rounding remain separate.
Stability with increasing $q$ cannot replace a bound, especially
near the phase singularities.

\newpage
\section{Retained failures and the next physical question}
The initial 22 scalar/algebra diagnostics contain two failures, preserved.
One used structural symbolic equality for a polynomial identity;
its exact difference simplifies to zero. The other tested a shell at
$\hbar=.5$, $\rho=.025$ and missed the original one-percent tolerance:
the deviation was $0.01009160698$.
A separate refinement keeps that tolerance and halves the radius;
the deviation becomes $0.00251886700$, with 64/128 angular samples
agreeing to the recorded precision. All five refinement checks pass.

A further 31 controls test the phase-vortex coefficient and local winding
using 50-digit products, two radii, and 32/64 angular samples.
They pass, but the quadrature controls are empirical.
Domain conclusions rely on the Laurent expansion and the
Lebesgue-point argument, not numerical extrapolation.

\paragraph{Next decisions.}
Compare any proposed weak or contour continuation with the bounded
paired functional calculus; specify its boundary values, pairing, and
test space. If they are proved equal, use that theorem directly.
If not, identify the additional assumption. A separate numerical route
is singularity-aware certified matrix-element quadrature,
including tails and rounding; more node counts alone do not provide it.

Even a resolved local pentagon leaves the observer problem:
physical insertion, mass/central-label correspondence, reflection,
and gluing measure must come from the same inner product before a
four-point Gram test is meaningful. A regular-pair representation on
discrete radial spectra changes the mathematical input, and cannot
be presented as the original de Sitter observer theory without that
dictionary. The complete holographic objective remains open.

\begin{thebibliography}{9}\small
\bibitem{gon} A. B. Goncharov,
\emph{Pentagon relation for the quantum dilogarithm and quantized
$M_{0,5}^{\rm cyc}$}, \href{https://arxiv.org/abs/0706.4054v2}{arXiv:0706.4054v2}.
\bibitem{fg} V. V. Fock and A. B. Goncharov,
\emph{The quantum dilogarithm and representations of quantum cluster varieties},
\href{https://arxiv.org/abs/math/0702397v6}{arXiv:math/0702397v6}.
\bibitem{ks} H. K. Kim and C. Scarinci,
\emph{A quantization of moduli spaces of 3-dimensional gravity},
\href{https://arxiv.org/abs/2112.13329v2}{arXiv:2112.13329v2}.
\bibitem{kim} H. K. Kim,
\emph{A trilogy of mapping class group representations from
three-dimensional quantum gravity},
\href{https://arxiv.org/abs/2308.02934v2}{arXiv:2308.02934v2}.
\bibitem{pseudo} H. K. Kim,
\emph{Three-dimensional quantum gravity from the quantum pseudo-K\"ahler plane},
\href{https://arxiv.org/abs/2112.13962v2}{arXiv:2112.13962v2}.
\bibitem{ps} W. Pusz and P. M. So{\l}tan,
\emph{Functional form of unitary representations of the quantum az+b group},
\href{https://arxiv.org/abs/math/0604613v1}{arXiv:math/0604613v1}.
\end{thebibliography}
\end{document}